\documentclass[11pt]{article}

\usepackage{amsmath,amssymb,amsthm}
\usepackage{geometry}
\geometry{margin=1in}

\title{Eine $\alpha$-kalibrierte Attraktorstruktur im diskreten Phasenraum\\
und die Stabilität asymmetrischer Lösungen}
\author{}
\date{}

\begin{document}
\maketitle

\begin{abstract}
Wir untersuchen einen diskreten, modular strukturierten Phasenraum, dessen Dynamik
durch eine $\alpha$-kalibrierte Kopplung bestimmt ist. Der Raum zerfällt in endlich
viele Äquivalenzklassen, die durch kohärente Resonanzbedingungen ausgezeichnet sind.
Es wird gezeigt, dass die Dynamik generisch zu einer asymmetrischen Attraktorstruktur
führt: kohärente Zustände besitzen offene Einzugsgebiete positiver Maßzahl, während
ihre symmetrischen Gegenorientierungen lediglich instabile Separatrizen bilden.
Die resultierende Asymmetrie ist klein, jedoch strukturell stabil und skaliert
universell mit der $\alpha$-Kalibrierung. Die Arbeit liefert eine rein mathematische
Erklärung für persistente Asymmetrien ohne explizite Symmetriebrechung auf Ebene der
Erzeugungsregeln.
\end{abstract}

\section{Einleitung}

Asymmetrische Langzeitlösungen in formal symmetrischen Dynamiken sind ein bekanntes
Phänomen nichtlinearer Systeme. Typischerweise werden solche Asymmetrien entweder durch
explizite Symmetriebrechung oder durch spezielle Anfangsbedingungen erklärt. In dieser
Arbeit wird ein dritter Mechanismus untersucht: asymmetrische Attraktorstruktur bei
symmetrischer Dynamik, hervorgerufen durch geometrisch-arithmetische Kalibrierung des
Phasenraums.

\section{Konstruktion des Phasenraums}

\subsection{Diskrete Zustandsmenge}

Sei $\mathcal{N}=\mathbb{N}$ die Menge der natürlichen Zahlen. Wir definieren eine
Äquivalenzrelation
\[
n \sim m \iff n \equiv m \pmod{12}
\]
und betrachten die vier Klassen
\[
\mathcal{F}=\{1,5,7,11\}\subset\mathbb{Z}_{12}.
\]

\subsection{Phasenabbildung}

Jedem Zustand $n\in\mathcal{N}$ wird eine Phase
\[
\theta(n)=2\pi\frac{\log n \bmod \log M}{\log M}
\]
zugeordnet. Der Phasenraum ist
\[
\mathcal{P}=S^1\times\mathcal{F}.
\]

\section{$\alpha$-kalibrierte Kopplung}

Wir definieren eine Kopplungsfunktion
\[
K_\alpha((\theta_i,f_i),(\theta_j,f_j))
=\alpha\cos(\theta_i-\theta_j)\chi(f_i,f_j),
\]
wobei $\chi$ symmetrisch und nichtnegativ ist.

\section{Attraktoren und Separatrizen}

Kohärente Attraktoren besitzen offene Einzugsgebiete positiver Maßzahl, während
gegenorientierte Zustände auf Separatrizen liegen und mindestens einen positiven
Lyapunov-Exponenten aufweisen.

\section{Maßtheoretische Asymmetrie}

Das Maß der Einzugsgebiete kohärenter Attraktoren ist strikt positiv, das der
gegenorientierten Zustände verschwindet.

\section{Skalierung}

Die relative Asymmetrie skaliert wie
\[
\Delta\sim\alpha^k,\qquad k\in\mathbb{N}.
\]

\section{Schlussfolgerung}

Die Asymmetrie entsteht ausschließlich aus der Attraktorstruktur eines formal
symmetrischen Systems.
\end{document}
